% संपूर्ण मराठी ग्रंथातील प्रकरण 52: प्रस्तावना
% हा संपादनयोग्य स्रोतखंड आहे. संकलनासाठी संपूर्ण स्रोतसंग्रहातील
% अक्षररूपे, पूर्वभाग, चिन्हव्याख्या आणि आकृती-स्रोत आवश्यक आहेत.
% मूळ: Open Logic Project; मजकूर आणि रूपांतर CC BY 4.0.
% यंत्रानुवाद, दुरुस्ती आणि तपासणी: OpenAI Codex —
% GPT-5.6 Sol आणि GPT-6 Sol, दोन्ही Ultra effort.
% दुरुस्ती-पुनरावलोकन: GPT-6.1 Sol, Ultra effort.
\chapter{प्रस्तावना}\label{int:int:chap}










\section{रचनाशील तर्कविचार}\label{int:int:cr:sec}

मोडल संकारक किंवा द्वितीय-क्रम संख्यापक जोडून अभिजात
तर्कशास्त्राचा विस्तार करता येतो; त्याउलट
अंतःप्रज्ञावादी तर्कशास्त्र अभिजात तर्कशास्त्रावर
मर्यादा घालते, म्हणून ते ``अभिजातेतर'' आहे.
अभिजात तर्कशास्त्र अनेक अर्थांनी \emph{अरचनाशील}
आहे. एका प्रकारच्या रचनाशील गणिताला वैशिष्ट्यपूर्ण
असणारा अधिक ``रचनाशील'' तर्कविचार अंतःप्रज्ञावादी
तर्कशास्त्रात पकडायचा आहे. यामागील काही प्रेरणा
पुढील उदाहरणांतून स्पष्ट होतील.

समजा, एखादी व्यक्ती असा दावा करते की तिने एक नैसर्गिक
संख्या~$n$ शोधली आहे: $n$ सम असेल, तर रीमान
गृहीतक सत्य आहे; आणि $n$ विषम असेल, तर ते असत्य
आहे. छान बातमी!{} रीमान गृहीतक सत्य आहे की नाही
हा गणितातील मोठ्या अनिर्णीत प्रश्नांपैकी एक आहे.
या व्यक्तीने तो प्रश्न गणनेचा, म्हणजे विशिष्ट संख्या
सम आहे की नाही हे ठरवण्याचा, प्रश्न केला असे दिसते.

मग~$n$ चे ते जादुई मूल्य कोणते? ती व्यक्ती असे
सांगते: रीमान गृहीतक सत्य असेल, तर~$n$ ही नैसर्गिक
संख्या $2$ आहे; अन्यथा ती $3$ आहे.

तुम्ही रागाने तुमचे पैसे परत मागता. अभिजात दृष्टिकोनातून
वरील वर्णन खरोखरच~$n$ चे एकमेव मूल्य ठरवते;
पण तुम्हाला हवे आहे ते~$n$ चे \emph{स्पष्टपणे}
दिलेले मूल्य.

आता आणखी एक, कदाचित कमी कृत्रिम, उदाहरण पाहू.
अपरिमेय संख्येचा परिमेय घात घेऊन परिमेय उत्तर मिळू
शकते, हे माहीत आहे. उदाहरणार्थ, $\sqrt{2}^2 = 2$.
पण अपरिमेय संख्येचा \emph{अपरिमेय} घात घेतल्यावरही
परिमेय उत्तर मिळू शकते का, हे तितके स्पष्ट नाही.
पुढील प्रमेय या प्रश्नाचे होकारार्थी उत्तर देते:

\begin{thm}
$a$ आणि $b$ अशा अपरिमेय संख्या आहेत की $a^b$ परिमेय आहे.
\end{thm}

\begin{proof}
$\sqrt{2}^{\sqrt{2}}$ पाहा. ही संख्या परिमेय असेल,
तर काम झाले: $a = b = \sqrt{2}$ घेता येईल. अन्यथा
ती अपरिमेय आहे. मग
\[(\sqrt{2}^{\sqrt{2}})^{\sqrt{2}} = \sqrt{2}^{\sqrt{2} \cdot
  \sqrt{2}} = \sqrt{2}^2 = 2,
\]
आणि हे परिमेय आहे. म्हणून या वेळी $a$ म्हणून
$\sqrt{2}^{\sqrt{2}}$ आणि $b$ म्हणून~$\sqrt 2$ घ्या.
\end{proof}

ही सिद्धता वैध आहे का? बहुतेक गणितज्ञांना ती वैध
वाटते. तरीही इथे थोडे असमाधानकारक काहीतरी आहे:
विशिष्ट गुणधर्म असलेल्या वास्तव संख्यांच्या जोडीचे
अस्तित्व सिद्ध केले, पण ती जोडी \emph{कोणती} हे
सांगता आले नाही. हाच निष्कर्ष अशा रीतीनेही सिद्ध
करता येतो की $a$, $b$ ही जोडी सिद्धतेत \emph{दिली}
जाते: $a = \sqrt{3}$ आणि $b = \log_3 4$ घ्या. मग
\[a^b = \sqrt{3}^{\log_3 4} = 3^{1/2 \cdot \log_3 4} = (3^{\log_3
  4})^{1/2} = 4^{1/2}= 2,
\]
कारण $3^{\log_3 x} = x$.

पहिल्या सिद्धतेतील पावलांसारखी पावले ग्राह्य नसतील
असा तर्कविचार अंतःप्रज्ञावादी तर्कशास्त्रात पकडायचा
आहे. $\varphi(x)$ ची पूर्तता करणारा $x$ अस्तित्वात आहे
हे सिद्ध करण्यासाठी दुसऱ्या सिद्धतेप्रमाणे एखादा
विशिष्ट~$x$ आणि तो $\varphi$ ची पूर्तता करतो याची
सिद्धता द्यावी लागते. $\varphi$ किंवा $\psi$ सत्य आहे
हे सिद्ध करायचे असेल, तर या दोहोंपैकी एकाची
सिद्धता देता आली पाहिजे.

औपचारिक भाषेत सांगायचे तर, अभिजात तर्कशास्त्राची
निष्पत्ती व्यवस्था विशिष्ट प्रकारे मर्यादित
केल्यावर अंतःप्रज्ञावादी तर्कशास्त्र मिळते. गणिती
दृष्टिकोनातून या केवळ औपचारिक निगमन व्यवस्था
आहेत; पण आधी सांगितल्याप्रमाणे त्यांत एका प्रकारचा
गणिती तर्कविचार पकडायचा आहे. तो तर्कविचार
गणिताविषयीच्या विशिष्ट तात्त्विक दृष्टिकोनातून
(उदाहरणार्थ, ब्रौवर यांच्या अंतःप्रज्ञावादातून)
समर्थनीय मानता येईल; किंवा बिशप यांच्या
रचनाशीलतावादाच्या धर्तीवर अधिक ``मूर्त'' आणि
समाधानकारक गणिती तर्कविचार मानता येईल. औपचारिक
वर्णन अनौपचारिक प्रेरणा खरोखर पकडते का, यावरही
वाद करता येईल. पण आपली तात्त्विक भूमिका काहीही
असो, अंतःप्रज्ञावादी तर्कशास्त्राचा औपचारिकरीत्या
मांडलेल्या तर्कशास्त्राप्रमाणे अभ्यास करता येतो;
आणि अनेक गणिती तर्कशास्त्रज्ञांना तो अभ्यास
रसपूर्ण वाटतो.












\section{अंतःप्रज्ञावादी तर्कशास्त्राची विन्यासमीमांसा}\label{int:int:syn:sec}

अंतःप्रज्ञावादी तर्कशास्त्राचा विन्यास विधानीय
तर्कशास्त्राच्या विन्यासासारखाच आहे. अभिजात विधानीय
तर्कशास्त्रात काही संयोजकांची व्याख्या इतर
संयोजकांनी करता येते. उदाहरणार्थ, $\varphi \lif \psi$ ची
$\lnot \varphi \lor \psi$ ने, किंवा $\varphi \lor \psi$ ची
$\lnot(\lnot \varphi \land \lnot \psi)$ ने व्याख्या करता
येते. म्हणून अभिजात तर्कशास्त्राच्या मांडणीत काही
संयोजक अशा व्याख्यांची संक्षिप्त रूपे म्हणून येतात.
अंतःप्रज्ञावादी तर्कशास्त्रात मात्र दोन अपवाद सोडता
तसे होत नाही: $\lnot \varphi$ हे $\varphi \lif \lfalse$ चे
संक्षिप्त रूप म्हणून परिभाषित करता येते---आणि
सहसा तसे करतात. अर्थात मग $\lfalse$ ची स्वतःची
व्याख्या इतर संयोजकांच्या साहाय्याने केलेली नसली पाहिजे!{} तसेच
$\varphi \liff \psi$ ची अभिजात तर्कशास्त्राप्रमाणे
$(\varphi \lif \psi) \land (\psi \lif \varphi)$ अशी व्याख्या
करता येते.

विधानीय अंतःप्रज्ञावादी तर्कशास्त्रातील सूत्रे
ही \emph{विधानीय चले} आणि विधानीय
स्थिरांक~$\lfalse$ यांपासून \emph{तार्किक संयोजके}
वापरून तयार होतात. आपल्याकडे पुढील गोष्टी आहेत:

\begin{enumerate}
\item विधानीय चलांचा गणनीय अनंत संच~$\PVar$: $\Obj p_0$,
  $\Obj p_1$, \dots
  \item असत्यता दर्शवणारा विधानीय स्थिरांक~$\lfalse$.
\item तार्किक संयोजके: $\land$
  (संयोग), $\lor$ (विकल्पयोग), $\lif$ (सशर्त)
\item विरामचिन्हे: (, ), आणि स्वल्पविराम.
\end{enumerate}

\begin{defn}[सूत्र]
\label{int:int:syn:defn:formulas} विधानीय अंतःप्रज्ञावादी
तर्कशास्त्रातील \emph{सूत्रे} चा संच~$\Frm[L_0]$
पुढीलप्रमाणे विगमनाने परिभाषित करतो:
\begin{enumerate}
\item $\lfalse$ हे अण्विक सूत्र आहे.

\item प्रत्येक विधानीय चल~$\Obj p_i$ हे अण्विक
  सूत्र आहे.

\item $\varphi$ आणि $\psi$ ही सूत्रे असतील, तर $(\varphi \land
  \psi)$ हे सूत्र आहे.

\item $\varphi$ आणि $\psi$ ही सूत्रे असतील, तर $(\varphi \lor \psi)$
  हे सूत्र आहे.

\item $\varphi$ आणि $\psi$ ही सूत्रे असतील, तर $(\varphi \lif \psi)$
  हे सूत्र आहे.

\item याखेरीज दुसरे काहीही सूत्र नाही.
\end{enumerate}
\end{defn}

वर दिलेल्या मूलभूत संयोजकांखेरीज पुढील
\emph{परिभाषित} चिन्हेही वापरतो: $\lnot$ (नकार)
आणि $\liff$ (द्विशर्त). परिभाषित
संकारकांनी तयार झालेल्या सूत्रांचा अर्थ पुढीलप्रमाणे:
\begin{enumerate}
\item $\lnot \varphi$ हे $\varphi \lif \lfalse$ चे संक्षिप्त रूप आहे.

\item $\varphi \liff \psi$ हे $(\varphi \lif \psi) \land (\psi
  \lif \varphi)$ चे संक्षिप्त रूप आहे.
\end{enumerate}

$\lnot$ ला औपचारिकरीत्या संक्षिप्त रूप मानले असले,
तरी कधीकधी ते मूलभूत असल्यासारखे~$\lnot$ चे स्पष्ट नियम
आणि व्याख्यांतील खंड देऊ. याचा मुख्य उद्देश सरावाचे
प्रश्न मांडता यावेत हा आहे.












\section{ब्रौवर--हीटिंग--कोल्मोगोरोव अर्थनिर्धारण}\label{int:int:bhk:sec}

\begin{editorial}
  BHK अर्थनिर्धारण वापरून अंतःप्रज्ञावादी विधानांची
  वैधता सिद्ध करणाऱ्या सिद्धता गोंधळात टाकतात;
  त्यांचे अधिक चांगले स्पष्टीकरण द्यायला हवे.
\end{editorial}

अंतःप्रज्ञावादी संयोजकांचे एक अनौपचारिक रचनाशील
अर्थनिर्धारण आहे; त्याला बहुधा ब्रौवर--हीटिंग--कोल्मोगोरोव
अर्थनिर्धारण म्हणतात. त्यात ``रचना'' ही संकल्पना
वापरली जाते; तिला रचनाशील सिद्धता समजू शकता.
(औपचारिक निष्पत्ती व्यवस्थेतील
निष्पत्ती शी गल्लत होऊ नये म्हणून BHK
अर्थनिर्धारणात ``सिद्धता'' हा शब्द वापरत नाही.)
या सहज समजणाऱ्या संकल्पनेच्या आधारे BHK अर्थनिर्धारण
अंतःप्रज्ञावादी संयोजकांचे अर्थ स्पष्ट करते.

\begin{enumerate}
\item अण्विक विधानाची रचना कशाला म्हणायचे हे आपल्याला
  माहीत आहे असे गृहीत धरू.
\item $\varphi_1 \land \varphi_2$ ची रचना म्हणजे $\tuple{M_1, M_2}$
  ही जोडी; येथे $M_1$ ही~$\varphi_1$ ची आणि $M_2$ ही~$\varphi_2$
  ची रचना आहे.
\item $\varphi_1 \lor \varphi_2$ ची रचना म्हणजे $\tuple{s, M}$
  ही जोडी; यात $s$ हे~$1$ आणि $M$ ही~$\varphi_1$ ची रचना
  असते, किंवा $s$ हे~$2$ आणि $M$ ही~$\varphi_2$ ची रचना
  असते.
\item $\varphi \lif \psi$ ची रचना म्हणजे~$\varphi$ ची रचना
  घेऊन तिला~$\psi$ च्या रचनेत रूपांतरित करणारे फलन.
\item $\lfalse$ (असंगती) ची कोणतीही रचना नसते.
\item $\lnot \varphi$ हे $\varphi \lif \lfalse$ चे समानार्थी रूप
  म्हणून परिभाषित केले आहे. म्हणजे $\lnot \varphi$ ची रचना
  ही~$\varphi$ ची रचना घेऊन तिला~$\lfalse$ च्या रचनेत
  रूपांतरित करणारे फलन असते.
\end{enumerate}

\begin{ex}
उदाहरणार्थ $\lnot \lfalse$ घ्या. त्याची रचना हे
असे फलन असते जे~$\lfalse$ ची कोणतीही रचना इनपुट
म्हणून घेऊन~$\lfalse$ ची रचना आउटपुट म्हणून देते.
ओळख फलन~$\Id{}$ ही अशीच एक रचना आहे:
$\lfalse$ ची रचना~$M$ दिली, तर $\Id{}(M) = M$
ही~$\lfalse$ ची रचना देते.
\end{ex}

साधारणपणे $\lnot \varphi$ चा अर्थ ``$\varphi$ ची रचना अशक्य आहे'' असा होतो.

\begin{ex}
कोणत्याही विधान~$\varphi$ साठी $\varphi \lif \lnot \lnot \varphi$
सिद्ध करू. हेच सूत्र $\varphi \lif ((\varphi \lif \lfalse) \lif
\lfalse)$ आहे. याची रचना असे फलन~$f$ असली पाहिजे
की $\varphi$ ची रचना~$M$ दिल्यावर ते $(\varphi \lif \lfalse)
\lif \lfalse$ ची रचना~$f(M)$ देते. हे फलन~$f$
$(\varphi \lif \lfalse) \lif \lfalse$ ची रचना कशी
तयार करते ते पाहू: $\varphi \lif \lfalse$
ची रचना~$h$ इनपुट म्हणून घेऊन~$\lfalse$ ची रचना
आउटपुट म्हणून देणारे फलन $g$ परिभाषित करायचे आहे.
त्यासाठी $g$ ची व्याख्या अशी करा: इनपुट~$h$ हे
आधी मिळालेल्या~$\varphi$ च्या
रचनेवर~$M$ लावा. $h$ चे आउटपुट~$h(M)$ ही
$\lfalse$ ची रचना असल्याने, $M$ ही~$\varphi$ ची रचना
असेल तेव्हा $f(M)(h) = h(M)$ ही~$\lfalse$ ची रचना
असते.
\end{ex}

\begin{ex}
आता $\lnot(\varphi \land \lnot \varphi)$, म्हणजेच $(\varphi
\land (\varphi \lif \lfalse)) \lif \lfalse$, याची रचना
देऊ. ती असे फलन~$f$ आहे की $\varphi \land (\varphi \lif
\lfalse)$ ची रचना~$M$ इनपुट म्हणून दिल्यास ते
$\lfalse$ ची रचना देते. संयोग $\psi_1 \land \psi_2$
ची रचना म्हणजे $\tuple{N_1, N_2}$ ही जोडी; येथे
$N_1$ ही~$\psi_1$ ची आणि $N_2$ ही~$\psi_2$ ची
रचना आहे. $\psi_1 \land \psi_2$ च्या रचनेतून
अनुक्रमे $\psi_1$ आणि $\psi_2$ यांच्या रचना मिळवणारी
फलने $p_1$ आणि $p_2$ अशी परिभाषित करू:
\begin{align*}
  p_1(\tuple{N_1, N_2}) &= N_1\\
  p_2(\tuple{N_1, N_2}) &= N_2
\end{align*}
फलन~$f$ पुढीलप्रमाणे काम करते. प्रथम ते आपल्या
इनपुट~$M$ ला $p_1$ लावते; त्यामुळे~$\varphi$ ची रचना
मिळते. नंतर~$M$ ला $p_2$ लावते; त्यामुळे
$\varphi \lif \lfalse$ ची रचना मिळते. ही रचना म्हणजे
फलन~$p_2(M)$; त्याला~$\varphi$ ची रचना इनपुट म्हणून
दिल्यास ते~$\lfalse$ ची रचना देते. दुसऱ्या शब्दांत,
$p_2(M)$ हे $p_1(M)$ ला लावल्यास~$\lfalse$ ची
रचना मिळते. म्हणून $f(M) = p_2(M)(p_1(M))$ अशी
व्याख्या करता येते.
\end{ex}

\begin{ex}
आता $((\varphi \land \psi) \lif \chi) \lif (\varphi \lif
(\psi \lif \chi))$ ची रचना देऊ. म्हणजे, $(\varphi \land
\psi) \lif \chi$ ची रचना~$g$ घेऊन $(\varphi \lif (\psi \lif
\chi))$ ची रचना देणारे फलन $f$ हवे. रचना~$g$
ही स्वतः फलन आहे: $\varphi \land \psi$ च्या रचनांपासून
$\chi$ च्या रचनांकडे जाणारे. आणि आउटपुट~$f(g)$
हे फलन~$h_g$ आहे: $\varphi$ च्या रचनांपासून, $\psi$
च्या रचनांना~$\chi$ च्या रचनांत नेणाऱ्या फलनांकडे जाणारे.

ठीक आहे, हे गोंधळात टाकणारे आहे. इनपुट~$g$ साठी
$f$ चे आउटपुट ठरणारे फलन $h_g$ तयार करायचे आहे.
$h_g$ चा इनपुट ही~$\varphi$ ची रचना~$M$ आहे.
$h_g(M)$ चे आउटपुट हे~$\psi$ च्या रचनांपैकी~$N$
पासून~$\chi$ च्या रचनांकडे जाणारे फलन~$k_M$ असले
पाहिजे. $k_{g,M}(N) = g(\tuple{M, N})$ असे ठेवू.
$\tuple{M, N}$ ही~$\varphi \land \psi$ ची रचना आहे हे
आठवा. म्हणून $k_{g,M}$ ही $\psi \lif \chi$ ची रचना
आहे: ती~$\psi$ च्या रचना~$N$ यांना~$\chi$ च्या
रचनांत नेते. आता $h_g(M) = k_{g,M}$ असे ठेवू.
हे फलन~$\varphi$ च्या रचना~$M$ यांना $\psi
\lif \chi$ च्या रचना~$k_{g,M}$ यांत नेते. पुढे
$f(g) = h_g$ असे ठेवू. हे फलन $(\varphi \land \psi)
\lif \chi$ च्या रचना~$g$ यांना $\varphi \lif (\psi \lif \chi)$
च्या रचनांत नेते. हुश्श!{}
\end{ex}

$\varphi \lor \lnot \varphi$ या विधानाला विमध्य सिद्धांत
म्हणतात. काही विशिष्ट~$\varphi$ साठी (उदाहरणार्थ,
$\lfalse \lor \lnot \lfalse$) तो सिद्ध करता येतो;
पण सर्वसाधारणपणे नाही. कारण अंतःप्रज्ञावादी
विकल्पयोगासाठी दोन विकल्पांपैकी एकाची रचना लागते;
पण काही विधाने अशी आहेत की सध्या ती सिद्धही
करता येत नाहीत किंवा खोडताही येत नाहीत
(उदाहरणार्थ, गोल्डबाखचे अनुमान). मात्र विमध्य
सिद्धांत खोडताही येत नाही: म्हणजे $\lnot \lnot (\varphi
\lor \lnot \varphi)$ सत्य असते.

\begin{ex}
$\lnot \lnot (\varphi \lor \lnot \varphi)$ सिद्ध करण्यासाठी
$\lnot(\varphi \lor \lnot \varphi)$, म्हणजे $(\varphi
\lor (\varphi \lif \lfalse)) \lif \lfalse$, याची रचना
घेऊन तिला~$\lfalse$ च्या रचनेत रूपांतरित करणारे
फलन~$f$ हवे. दुसऱ्या शब्दांत, $g$ ही $\lnot(\varphi
\lor \lnot \varphi)$ ची रचना असेल, तर $f(g)$ ही
$\lfalse$ ची रचना ठरेल, असे फलन $f$ हवे.

$g$ ही $\lnot(\varphi \lor \lnot \varphi)$ ची रचना आहे असे
समजा; म्हणजे ते $\varphi \lor \lnot \varphi$ ची रचना घेऊन
तिला~$\lfalse$ च्या रचनेत नेणारे फलन आहे.
$\varphi \lor \lnot \varphi$ ची रचना म्हणजे $\tuple{s, M}$
ही जोडी: यात एकतर $s = 1$ आणि $M$ ही~$\varphi$ ची
रचना असते, किंवा $s = 2$ आणि $M$ ही $\lnot \varphi$
ची रचना असते. $h_1$ हे~$\varphi$ ची रचना~$M_1$ घेऊन
$\varphi \lor \lnot \varphi$ ची रचना देणारे फलन असू द्या:
ते $M_1$ ला $\tuple{1, M_1}$ कडे नेते. तसेच $h_2$
हे $\lnot \varphi$ ची रचना~$M_2$ घेऊन
$\varphi \lor \lnot \varphi$ ची रचना देणारे फलन असू द्या:
ते $M_2$ ला $\tuple{2, M_2}$ कडे नेते.

$k$ हे $\comp{h_1}{g}$ असू द्या. ते~$\varphi$ ची रचना
दिल्यावर~$\lfalse$ ची रचना देणारे फलन आहे; म्हणजे
ते~$\varphi \lif \lfalse$ किंवा $\lnot \varphi$ ची रचना
आहे. आता $l$ हे $\comp{h_2}{g}$ असू द्या. हे
फलन $\lnot \varphi$ ची रचना दिल्यावर~$\lfalse$ ची
रचना देते. $k$ ही $\lnot \varphi$ ची रचना असल्याने,
$l(k)$ ही~$\lfalse$ ची रचना आहे.

अशा प्रकारे $\lnot(\varphi \lor \lnot \varphi)$ ची
रचना~$g$ घेऊन तिला~$\lfalse$ च्या रचनेत कसे
रूपांतरित करायचे ते सांगितले. म्हणजे
$\lnot(\varphi \lor \lnot \varphi)$ ची रचना~$g$
घेऊन~$\lfalse$ ची रचना~$l(k)$ देणारे फलन $f$
ही $\lnot\lnot(\varphi \lor \lnot \varphi)$ ची रचना आहे.
\end{ex}

पाहिल्याप्रमाणे, सूत्रांची अंतःप्रज्ञावादी
वैधता दाखवण्यासाठी BHK अर्थनिर्धारण वापरणे लवकरच
अवजड आणि गोंधळात टाकणारे होते. सुदैवाने,
अंतःप्रज्ञावादी तर्कशास्त्रासाठी अधिक चांगल्या
निष्पत्ती व्यवस्था आणि अधिक नेमकी
चिन्हार्थमीमांसक अर्थलक्षणे उपलब्ध आहेत.











\section{नैसर्गिक निगमन}\label{int:int:ntd:sec}

$\FalseCl$ चे नियम वगळलेले नैसर्गिक निगमन ही
अंतःप्रज्ञावादी तर्कशास्त्रासाठीची प्रमाणित
निष्पत्ती व्यवस्था आहे. येथे ते नियम पुन्हा
देऊन BHK अर्थनिर्धारणाच्या साहाय्याने त्यामागची
प्रेरणा स्पष्ट करू. प्रत्येक नियम असा समजू शकतो
की आधारविधानांच्या रचना असतील, तर निष्कर्षाचीही
रचना आहे असे त्यावरून ठरवता येते.

नैसर्गिक निगमनातील निष्पत्त्यांमध्ये मुक्त न
केलेली गृहीतके असतात. म्हणून~$\Gamma$ या
मुक्त न केलेले गृहीतकांपासून~$\varphi$ ची
निष्पत्ती असेल, तर तिला प्रत्येक $\psi \in
\Gamma$ च्या रचना घेऊन~$\varphi$ ची रचना देणारे
फलन मानावे. कोणत्याही मुक्त न केलेले गृहीतकांशिवाय
$\varphi$ ची निष्पत्ती असेल, तर BHK अर्थनिर्धारणाच्या
अर्थाने~$\varphi$ ची रचना असते. मात्र पुढील चर्चेत गरज
नसेल तेव्हा~$\Gamma$ चा उल्लेख वगळू.

गृहीतक~$\varphi$ स्वतःच त्या मुक्त न केलेले
गृहीतक~$\varphi$ पासून~$\varphi$ ची निष्पत्ती आहे.
हे BHK अर्थनिर्धारणाशी जुळते: रचनांवरील ओळख फलन
$\varphi$ ची कोणतीही रचना घेऊन~$\varphi$ चीच रचना देते.

\subsection{संयोग}

\begin{defish}
\AxiomC{$\varphi$}
\AxiomC{$\psi$}
\RightLabel{\Intro{\land}}
\BinaryInfC{$\varphi \land \psi$}
\DisplayProof
\hfill
\begin{tabular}{l}
\AxiomC{$\varphi \land \psi$}
\RightLabel{\Elim{\land}}
\UnaryInfC{$\varphi$}
\DisplayProof
\\[3ex]
\AxiomC{$\varphi \land \psi$}
\RightLabel{\Elim{\land}}
\UnaryInfC{$\psi$}
\DisplayProof
\end{tabular}
\end{defish}
आपल्याकडे अनुक्रमे $\varphi_1$ आणि $\varphi_2$ यांच्या
रचना $N_1$, $N_2$ आहेत असे समजा. मग
$\varphi_1 \land \varphi_2$ चीही रचना आहे: जोडी
$\tuple{N_1, N_2}$.
BHK अर्थनिर्धारणानुसार $\varphi_1 \land \varphi_2$ ची
रचना म्हणजे $\tuple{N_1, N_2}$ ही जोडी. अशी जोडी
आपल्याकडे आहे असे समजा. मग प्रत्येक संयुक्त घटकाची
रचना आहे: $N_1$ ही~$\varphi_1$ ची आणि $N_2$ ही
$\varphi_2$ ची रचना आहे.
\subsection{सशर्त}
\begin{defish}
  \AxiomC{$\Discharge{\varphi}{u}$}
  \DeduceC{$\psi$}
  \DischargeRule{$\Intro{\lif}$}{u}
  \UnaryInfC{$\varphi \lif \psi$}
  \DisplayProof
\hfill
  \AxiomC{$\varphi \lif \psi$}
  \AxiomC{$\varphi$}
  \RightLabel{$\Elim{\lif}$}
  \BinaryInfC{$\psi$}
  \DisplayProof
\end{defish}
मुक्त न केलेले गृहीतक~$\varphi$ पासून~$\psi$ ची
निष्पत्ती असेल, तर~$\varphi$ च्या रचनांना~$\psi$
च्या रचनांत नेणारे फलन~$f$ असते. तेच फलन
$\varphi \lif \psi$ ची रचना आहे. म्हणून $\Intro\lif$
च्या आधारविधानाची रचना~$\varphi$ च्या रचनेवर अवलंबून
असेल, तर निष्कर्ष $\varphi
\lif \psi$ ची रचना असते.
दुसरीकडे, $\varphi$ ची रचना~$N$ आणि $\varphi \lif \psi$
ची रचना~$f$ आहेत असे समजा. $\varphi \lif \psi$ ची
रचना हे~$\varphi$ च्या रचनांना~$\psi$ च्या रचनांत नेणारे
फलन असते. म्हणून $f(N)$ ही~$\psi$ ची रचना आहे;
म्हणजे $\Elim\lif$ च्या निष्कर्षाची रचना आहे.
\subsection{विकल्पयोग}
\begin{defish}
\begin{tabular}{l}
\AxiomC{$\varphi$}
\RightLabel{\Intro{\lor}}
\UnaryInfC{$\varphi \lor \psi$}
\DisplayProof
\\[3ex]
\AxiomC{$\psi$}
\RightLabel{\Intro{\lor}}
\UnaryInfC{$\varphi \lor \psi$}
\DisplayProof
\end{tabular}
\hfill
\AxiomC{$\varphi \lor \psi$}
\AxiomC{$\Discharge{\varphi}{n}$}
\DeduceC{$\chi$}
\AxiomC{$\Discharge{\psi}{n}$}
\DeduceC{$\chi$}
\DischargeRule{\Elim{\lor}}{n}
\TrinaryInfC{$\chi$}
\DisplayProof
\end{defish}
$\varphi_i$ ची रचना~$N_i$ असेल, तर तिच्यापासून
$\varphi_1 \lor \varphi_2$ ची रचना~$\tuple{i, N_i}$ करता
येते. उलट, $\varphi_1 \lor \varphi_2$ ची रचना म्हणजे
$\tuple{i, N_i}$ ही जोडी आपल्याकडे असेल, जिथे
$N_i$ ही~$\varphi_i$ ची रचना आहे; तसेच अनुक्रमे
$\varphi_1$, $\varphi_2$ यांच्या रचनांना~$\chi$ च्या रचनांत
नेणारी फलने $f_1$, $f_2$ असतील, तर $f_i(N_i)$
ही~$\chi$ ची रचना मिळते. हाच~$\Elim\lor$ चा निष्कर्ष.
\subsection{असंगती}
\begin{defish}
  \AxiomC{$\lfalse$}
  \RightLabel{\FalseInt}
  \UnaryInfC{$\varphi$}
  \DisplayProof
\end{defish}
मुक्त न केलेले गृहीतके~$\psi_1$, \dots, $\psi_n$
यांपासून~$\lfalse$ ची निष्पत्ती असेल, तर
$\psi_1$, \dots, $\psi_n$ यांच्या रचनांना~$\lfalse$
च्या रचनेत नेणारे फलन~$f(M_1,
\dots, M_n)$ असते. $\lfalse$ ची रचना नसल्याने
$\psi_1$, \dots, $\psi_n$ या सर्वांच्या रचनाही एकत्र
असू शकत नाहीत. म्हणून $f$ बाबत पुढीलही खरे असते:
$M_1$, \dots, $M_n$ या अनुक्रमे $\psi_1$, \dots,
$\psi_n$ यांच्या रचना \emph{असतील तर}, $f(M_1, \dots,
M_n)$ ही~$\varphi$ ची रचना \emph{असते}.
\subsection{$\lnot$ साठीचे नियम}
$\lnot \varphi$ ची व्याख्या $\varphi \lif \lfalse$ अशी
असल्याने काटेकोरपणे पाहता~$\lnot$ साठी वेगळे
नियम लागत नाहीत. तरी ते द्यायचे झाले, तर असे दिसतील:
\begin{defish}
\AxiomC{$\Discharge{\varphi}{n}$}
\noLine
\DeduceC{$\lfalse$}
\DischargeRule{\Intro{\lnot}}{n}
\UnaryInfC{$\lnot \varphi$}
\DisplayProof
\hfill
\AxiomC{$\lnot \varphi$}
\AxiomC{$\varphi$}
\RightLabel{\Elim{\lnot}}
\BinaryInfC{$\lfalse$}
\DisplayProof
\end{defish}
\subsection{निष्पत्तीची उदाहरणे}
\begin{enumerate}
\item $\Proves \varphi \lif (\lnot \varphi \lif \lfalse)$, i.e., $\Proves \varphi
  \lif ((\varphi \lif \lfalse) \lif \lfalse)$
  \begin{prooftree}
    \AxiomC{$\Discharge{\varphi}{2}$}
    \AxiomC{$\Discharge{\varphi \lif \lfalse}{1}$}
    \RightLabel{\Elim\lif}
    \BinaryInfC{$\lfalse$}
    \DischargeRule{\Intro\lif}{1}
    \UnaryInfC{$(\varphi \lif \lfalse)\lif \lfalse$}
    \DischargeRule{\Intro\lif}{2}
    \UnaryInfC{$\varphi \lif (\varphi \lif \lfalse) \lif \lfalse$}
  \end{prooftree}
\item $\Proves ((\varphi \land \psi) \lif \chi) \lif (\varphi \lif (\psi \lif \chi))$
  \begin{prooftree}
    \AxiomC{$\Discharge{(\varphi \land \psi) \lif \chi}{3}$}
    \AxiomC{$\Discharge{\varphi}{2}$}
    \AxiomC{$\Discharge{\psi}{1}$}
    \RightLabel{\Intro\land}
    \BinaryInfC{$\varphi \land \psi$}
    \RightLabel{\Elim\lif}
    \BinaryInfC{$\chi$}
    \DischargeRule{\Intro\lif}{1}
    \UnaryInfC{$\psi \lif \chi$}
    \DischargeRule{\Intro\lif}{2}
    \UnaryInfC{$\varphi \lif (\psi \lif \chi)$}
    \DischargeRule{\Intro\lif}{3}
    \UnaryInfC{$((\varphi \land \psi) \lif \chi) \lif (\varphi \lif (\psi \lif \chi))$}
  \end{prooftree}
\item $\Proves \lnot(\varphi \land \lnot \varphi)$, i.e., $\Proves (\varphi \land (\varphi
  \lif \lfalse))\lif \lfalse$
  \begin{prooftree}
    \AxiomC{$\Discharge{\varphi \land (\varphi \lif \lfalse)}{1}$}
    \RightLabel{\Elim\land}
    \UnaryInfC{$\varphi \lif \lfalse$}
    \AxiomC{$\Discharge{\varphi \land (\varphi \lif \lfalse)}{1}$}
    \RightLabel{\Elim\land}
    \UnaryInfC{$\varphi$}
    \RightLabel{\Elim\lif}
    \BinaryInfC{$\lfalse$}
    \DischargeRule{\Intro\lif}{1}
    \UnaryInfC{$(\varphi \land (\varphi \lif \lfalse))\lif \lfalse$}
  \end{prooftree}
\item $\Proves \lnot\lnot(\varphi \lor \lnot \varphi)$, i.e., $\Proves ((\varphi \lor
  (\varphi \lif \lfalse))\lif \lfalse)\lif \lfalse$
  \begin{prooftree}\hskip-3em
    \AxiomC{$\Discharge{(\varphi \lor (\varphi \lif \lfalse))\lif \lfalse}{2}$}
    \AxiomC{$\Discharge{(\varphi \lor (\varphi \lif \lfalse))\lif \lfalse}{2}$}
    \AxiomC{$\Discharge{\varphi}{1}$}
    \RightLabel{\Intro\lor}
    \UnaryInfC{$\varphi \lor (\varphi \lif \lfalse)$}
    \RightLabel{\Elim\lif}
    \BinaryInfC{$\lfalse$}
    \DischargeRule{\Intro\lif}{1}
    \UnaryInfC{$\varphi \lif \lfalse$}
    \RightLabel{\Intro\lor}
    \UnaryInfC{$\varphi \lor (\varphi \lif \lfalse)$}
    \RightLabel{\Elim\lif}
    \insertBetweenHyps{\hskip -4em}
    \BinaryInfC{$\lfalse$}
    \DischargeRule{\Intro\lif}{2}
    \UnaryInfC{$((\varphi \lor (\varphi \lif \lfalse))\lif \lfalse)\lif \lfalse$}
  \end{prooftree}
\end{enumerate}
\begin{prop}
  अंतःप्रज्ञावादी तर्कशास्त्रातील नैसर्गिक निगमनात
  $\Gamma \Proves \varphi$ असेल, तर अभिजात तर्कशास्त्रातही
  $\Gamma \Proves \varphi$ असते. विशेषतः, $\varphi$ हे
  अंतःप्रज्ञावादी प्रमेय असेल, तर ते अभिजात प्रमेयही असते.
\end{prop}
\begin{proof}
  नैसर्गिक निगमनाचा प्रत्येक नियम अभिजात नैसर्गिक
  निगमनाचाही नियम आहे. म्हणून अंतःप्रज्ञावादी
  तर्कशास्त्रातील प्रत्येक निष्पत्ती ही अभिजात
  तर्कशास्त्रातीलही निष्पत्ती आहे.
\end{proof}
\begin{prob}
  पुढील सूत्रांची अंतःप्रज्ञावादी तर्कशास्त्रात
  निष्पत्त्या द्या:
  \begin{enumerate}
    \item $(\lnot \varphi \lor \psi) \lif (\varphi \lif \psi)$
    \item $\lnot\lnot\lnot \varphi \lif \lnot \varphi$
    \item $\lnot\lnot (\varphi \land \psi) \liff (\lnot\lnot \varphi\land \lnot \lnot \psi)$
    \item $\lnot (\varphi \lor \psi) \liff (\lnot \varphi \land \lnot \psi)$
    \item $(\lnot \varphi \lor \lnot \psi) \lif \lnot (\varphi \land \psi)$
    \item $\lnot\lnot ( \varphi \land \psi) \lif  (\lnot\lnot \varphi \lor
    \lnot\lnot \psi)$
  \end{enumerate}
\end{prob}
\section{स्वयंसिद्धकीय निष्पत्ती}\label{int:int:axd:sec}
अंतःप्रज्ञावादी विधानीय तर्कशास्त्रातील स्वयंसिद्धकीय
निष्पत्त्या या संकल्पनात्मकदृष्ट्या सर्वांत सोप्या
आणि ऐतिहासिकदृष्ट्या पहिल्या निष्पत्ती-पद्धती आहेत. त्यांचे कार्य अभिजात विधानीय
तर्कशास्त्रातल्याप्रमाणेच चालते.
\begin{defn}[निष्पन्नता]
जर $\Gamma$ हा $\Lang L$ मधील सूत्रांचा संच
असेल, तर $\Gamma$ पासूनची \emph{निष्पत्ती} म्हणजे
सूत्रांचा सांत अनुक्रम $\varphi_1$, \dots,~$\varphi_n$
असा की प्रत्येक $i \le n$ साठी पुढीलपैकी एक अट
पूर्ण होते:
\begin{enumerate}
\item $\varphi_i \in \Gamma$; किंवा
\item $\varphi_i$ हे स्वयंसिद्धक आहे; किंवा
\item एखाद्या $\varphi_j$ आणि $\varphi_k$ पासून, जिथे $j < i$
  आणि $k < i$, विध्यात्मक अनुमानाने $\varphi_i$ निष्पन्न होते;
  म्हणजे $\varphi_k \ident \varphi_j \lif \varphi_i$.
\end{enumerate}
\end{defn}
\begin{defn}[स्वयंसिद्धके]
अंतःप्रज्ञावादी विधानीय तर्कशास्त्रातील
\emph{स्वयंसिद्धकांचा} संच $\PAx$ हा पुढील रूपांच्या
सर्व सूत्रांचा बनलेला आहे:
\begin{align}
  & (\varphi \land \psi) \lif \varphi \label{int:int:axd:ax:land1}\\
  & (\varphi \land \psi) \lif \psi \label{int:int:axd:ax:land2}\\
  & \varphi \lif (\psi \lif (\varphi \land \psi)) \label{int:int:axd:ax:land3}\\
  & \varphi \lif (\varphi \lor \psi) \label{int:int:axd:ax:lor1}\\
  & \varphi \lif (\psi \lor \varphi) \label{int:int:axd:ax:lor2}\\
  & (\varphi \lif \chi) \lif ((\psi \lif \chi) \lif ((\varphi \lor \psi) \lif \chi)) \label{int:int:axd:ax:lor3}\\
  & \varphi \lif (\psi \lif \varphi) \label{int:int:axd:ax:lif1}\\
  & (\varphi \lif (\psi \lif \chi)) \lif ((\varphi \lif \psi) \lif (\varphi \lif \chi)) \label{int:int:axd:ax:lif2}\\
  & \lfalse \lif \varphi \label{int:int:axd:ax:lfalse1}
\end{align}
\end{defn}
\begin{defn}[निष्पन्नता]
सूत्र~$\varphi$ हे $\Gamma$ पासून \emph{निष्पन्न करता येण्याजोगे}
असते, म्हणजे $\Gamma \Proves \varphi$, जर $\Gamma$
पासून सुरू होऊन $\varphi$ वर संपणारी निष्पत्ती असेल.
\end{defn}
\begin{defn}[प्रमेये]
रिक्त संचापासून $\varphi$ ची निष्पत्ती असेल,
तर सूत्र~$\varphi$ हे \emph{प्रमेय} असते.
$\varphi$ हे प्रमेय असेल तर आपण $\Proves \varphi$ लिहितो
आणि ते प्रमेय नसेल तर $\Proves/ \varphi$ लिहितो.
\end{defn}
\begin{prop}
  अंतःप्रज्ञावादी तर्कशास्त्राच्या स्वयंसिद्धकीय
  निष्पत्ती-पद्धतीत $\Gamma \Proves \varphi$
  असेल, तर अभिजात तर्कशास्त्रातही
  $\Gamma \Proves \varphi$ असते. विशेषतः, $\varphi$ हे
  अंतःप्रज्ञावादी प्रमेय असेल, तर ते अभिजात प्रमेयही असते.
\end{prop}
\begin{proof}
  प्रत्येक अंतःप्रज्ञावादी स्वयंसिद्धक हे अभिजात
  स्वयंसिद्धकही आहे. म्हणून अंतःप्रज्ञावादी
  तर्कशास्त्रातील प्रत्येक निष्पत्ती ही अभिजात
  तर्कशास्त्रातीलही निष्पत्ती आहे.
\end{proof}
